% PAGES: 331-334

Note that, by definition, $L_1, L_2, U_1$, and $U_2$ are nonsingular matrices; see
Definition 2.1. Multiply (10.143) by $L_2^{-1}$ on the left and by $U_1^{-1}$ on the
right and obtain
\begin{align}
L_2^{-1}\cdot(L_1U_1)\cdot U_1^{-1} &= L_2^{-1}\cdot(L_2U_2)\cdot U_1^{-1} \notag \\
\iff (L_2^{-1}L_1)\cdot(U_1U_1^{-1}) &= (L_2^{-1}L_2)\cdot(U_2U_1^{-1}) \notag \\
\iff L_2^{-1}L_1 &= U_2U_1^{-1},
\end{align}
since $U_1U_1^{-1} = I$ and $L_2^{-1}L_2 = I$.

Recall from Lemma 1.17 that the inverse of a lower triangular matrix is lower
triangular, and the inverse of an upper triangular matrix is upper triangular, and,
from Lemma 1.15, that the product of two lower triangular matrices is lower
triangular, and the product of two upper triangular matrices is upper triangular.
Thus, the matrix $L_2^{-1}$ is lower triangular and therefore the matrix
$L_2^{-1}L_1$ is also lower triangular. Similarly, the matrix $U_1^{-1}$ is upper
triangular and therefore the matrix $U_2U_1^{-1}$ is also upper triangular. Since
the matrices $L_2^{-1}L_1$ and $U_2U_1^{-1}$ are equal, see (10.144), it follows
that they must be diagonal matrices.

Let $D = \diag(d_k)_{k=1:n}$ be a diagonal matrix such that
\begin{equation}
D \;=\; L_2^{-1}L_1 \;=\; U_2U_1^{-1}.
\end{equation}

By multiplying $D = L_2^{-1}L_1$ to the left by $L_2$, we find that
\begin{equation}
L_2D \;=\; L_2\cdot(L_2^{-1}L_1) \;=\; (L_2^{-1}L_2)\cdot L_1 \;=\; L_1.
\end{equation}

Let $L_2 = \mathrm{col}\left(l_k^{(2)}\right)_{k=1:n}$. From Lemma 1.10, we obtain
that
\begin{equation}
L_2D \;=\; \mathrm{col}\left(d_kl_k^{(2)}\right)_{k=1:n}.
\end{equation}

Recall that all the diagonal entries of the matrices $L_1$ and $L_2$ are equal to
1. From (10.147), it follows that the $k$-th diagonal entry of the matrix $L_2D$ is
$d_k$, for all $k=1:n$. Since $L_2D = L_1$, see (10.146), and since all the diagonal
entries of $L_1$ are equal to 1, we obtain that $d_k=1$ for all $k=1:n$, and
therefore the matrix $D$ is the identity matrix, i.e., $D = I$.

Then, from (10.145), we find that
\[
I \;=\; L_2^{-1}L_1 \;=\; U_2U_1^{-1}.
\]

Thus, $L_1 = L_2$ and $U_1 = U_2$, and we conclude that the LU decomposition of the
matrix $A$ is unique.
\end{proof}

\section{More technical results}

We include here the proof of Lemma 1.17 stating that the inverse of a triangular
matrix is triangular of the same type:

\begin{lemma}
(i) The inverse of an upper triangular matrix is upper triangular.

(ii) The inverse of a lower triangular matrix is lower triangular.
\end{lemma}

\begin{proof}
(i) Let $U = \mathrm{col}(u_k)_{k=1:n}$ be a nonsingular upper triangular matrix and
let $B = \mathrm{col}(b_k)_{k=1:n}$ be the inverse matrix of $U$. Then,
\begin{equation}
BU \;=\; UB \;=\; I.
\end{equation}

From (1.11), we obtain that $BU = \mathrm{col}(Bu_k)_{k=1:n}$. Since $I =
\mathrm{col}(e_k)_{k=1:n}$, we find from (10.148) that
\begin{equation}
Bu_k \;=\; e_k, \quad \forall\, k = 1:n.
\end{equation}

Let $k=1$ in (10.149). Then, $Bu_1 = e_1$. Since $U$ is upper triangular, it
follows from (1.107) that $u_1(i) = 0$, for all $2 \le i \le n$. Thus,
\[
e_1 \;=\; Bu_1 \;=\; \sum_{i=1}^n u_1(i)b_i \;=\; u_1(1)b_1,
\]
and therefore
\[
b_1 \;=\; \frac{1}{u_1(1)}\,e_1.
\]

Note that $u_1(1) = U(1,1) \ne 0$ since $U$ is nonsingular; cf. Lemma 1.16.

We conclude that all the entries of the first column of the matrix $B$ with the
exception of the first entry are equal to 0, i.e., the first column of the matrix
$B$ satisfies the equivalent definition (1.107) of an upper triangular matrix.

We show by (complete) induction that all the columns of the matrix $B$ satisfy
property (1.107).

Let $k$ such that $2 \le k \le n$. Assume that the columns $b_1, \ldots, b_{k-1}$
satisfy property (1.107), i.e.,
\begin{equation}
b_i(j) \;=\; 0, \quad \forall\, i = 1:(k-1), \quad \forall\, i+1 \le j \le n.
\end{equation}
Since $U$ is upper triangular, we also find from (1.107) that
\begin{equation}
u_k(i) \;=\; 0, \quad \forall\, k+1 \le i \le n.
\end{equation}

Recall from (10.149) that $Bu_k = e_k$. From (1.7) and (10.151), it follows that
\begin{align}
e_k &= Bu_k \;=\; \sum_{i=1}^n u_k(i)b_i \;=\; \sum_{i=1}^k u_k(i)b_i \notag \\
&= u_k(k)b_k \;+\; \sum_{i=1}^{k-1} u_k(i)b_i,
\end{align}
and therefore
\begin{equation}
b_k \;=\; \frac{1}{u_k(k)}\left(e_k - \sum_{i=1}^{k-1} u_k(i)b_i\right).
\end{equation}

Note that $u_k(k) = U(k,k) \ne 0$ since $U$ is nonsingular; cf. Lemma 1.16.

Let $j$ such that $k+1 \le j \le n$. Then, $e_k(j) = 0$, and, for any $1 \le i \le
k-1$, $b_i(j) = 0$ as well, from the induction hypothesis (10.150).

Then, from (10.152), it follows that $b_k(j) = 0$ for any $k+1 \le j \le n$, and
therefore that the column $b_k$ also satisfies property (1.107).

By induction, we conclude that all the columns of the matrix $B$ satisfy property
(1.107), and therefore that $B$ is an upper triangular matrix.

(ii) Let $L$ be a nonsingular lower triangular matrix, and let $L^{-1}$ be the
inverse matrix of $L$. Then, by definition,
\begin{equation}
LL^{-1} \;=\; L^{-1}L \;=\; I.
\end{equation}

By transposing (10.153) and using (1.24), we obtain that
\[
(L^{-1})^tL^t \;=\; L^t(L^{-1})^t \;=\; I.
\]

Thus, $(L^{-1})^t$ is the inverse matrix of $L^t$, which is an upper triangular
matrix.

We already proved that the inverse of an upper triangular matrix is upper
triangular. It then follows that $(L^{-1})^t$ is an upper triangular matrix, and we
therefore conclude that $L^{-1}$ is a lower triangular matrix.
\end{proof}

The result below was used to establish (9.90) in section 9.6:

\begin{lemma}
If $Z$ is a standard normal variable, then
\[
P(Z \le -a) \;=\; 1 - P(Z \le a), \quad \forall\, a \in \R.
\]
\end{lemma}

\begin{proof}
Recall that
\begin{equation}
P(Z \le -a) \;=\; \frac{1}{\sqrt{2\pi}}\int_{-\infty}^{-a} e^{-\frac{x^2}{2}}\,dx,
\end{equation}
since the probability density function of the standard normal variable $Z$ is
$\frac{1}{\sqrt{2\pi}}e^{-\frac{x^2}{2}}$; see (7.109). Using the substitution $x =
-y$ in (10.154), we find that
\begin{align}
P(Z \le -a) &= \frac{1}{\sqrt{2\pi}}\int_\infty^a e^{-\frac{(-y)^2}{2}}(-dy)
\;=\; \frac{1}{\sqrt{2\pi}}\int_a^\infty e^{-\frac{y^2}{2}}\,dy \notag \\
&= 1 \;-\; \frac{1}{\sqrt{2\pi}}\int_{-\infty}^a e^{-\frac{y^2}{2}}\,dy \notag \\
&= 1 \;-\; P(Z \le a),
\end{align}
which is what we wanted to show. Note that, for the third equality above, we used
the fact that
\[
\frac{1}{\sqrt{2\pi}}\int_{-\infty}^a e^{-\frac{y^2}{2}}\,dy \;+\;
\frac{1}{\sqrt{2\pi}}\int_a^\infty e^{-\frac{y^2}{2}}\,dy \;=\;
\frac{1}{\sqrt{2\pi}}\int_{-\infty}^\infty e^{-\frac{y^2}{2}}\,dy \;=\; 1,
\]
since the integral over $(-\infty,\infty)$ of the probability density
$\frac{1}{\sqrt{2\pi}}e^{-\frac{y^2}{2}}$ of $Z$ is equal to 1.
\end{proof}

\section{Exercises}

\begin{exercises}

\item Show that
\[
\det\begin{pmatrix} 1 & a & a^2 \\ 1 & b & b^2 \\ 1 & c & c^2 \end{pmatrix}
\;=\; (c-a)(c-b)(b-a),
\]
where $a,b,c \in \R$.

Note: A matrix of the form
\[
\begin{pmatrix} 1 & c_1 & c_1^2 & \ldots & c_1^{n-1} \\
1 & c_2 & c_2^2 & \ldots & c_2^{n-1} \\
\vdots & \vdots & \vdots & \ddots & \vdots \\
1 & c_n & c_n^2 & \ldots & c_n^{n-1} \end{pmatrix},
\]
where $c_1, c_2, \ldots, c_n$ are constants, is called a Vandermonde matrix. The
determinant of the Vandermonde matrix is equal to
\[
\prod_{1 \le k < j \le n} (c_j - c_k).
\]

\item Show that any orthogonal matrix has determinant 1 or $-1$. In other words,
show that, for any orthogonal matrix $Q$, either $\det(Q) = 1$, or $\det(Q) = -1$.

\item Let $M_1$ and $M_2$ be $n \times n$ symmetric matrices. Show that, if
\[
x^tM_1x \;=\; x^tM_2x, \quad \forall\, x \in \R^n,
\]
then $M_1 = M_2$.

\item Show that the quadratic form of a matrix is 0 if and only if the matrix is
skew--symmetric, i.e., show that
\[
q_A(x) = 0 \quad \text{if and only if} \quad A^t = -A.
\]

\end{exercises}

% This is the last chunk of Chapter 10 (the book's final chapter). PDF page 334
% (folio 318) is the true end of the chapter: after Exercise 4 the page is blank
% down to the bottom margin, and PDF page 335 opens the book's bibliography
% (out of scope for this chunk). No environment, display math, or non-environment
% example/answer/remark block is left open at the end of this file.
